\documentclass[11pt]{article}
\usepackage[a4paper,margin=1in]{geometry}
\usepackage{amsmath,amssymb,mathtools}
\usepackage[hidelinks]{hyperref}
\setlength{\parindent}{0pt}
\setlength{\parskip}{0.6em}

\title{Status of Audenaert--Kittaneh Problem 4\\
\large An explicit 2014 disproof and a supplementary exact $2\times2$ witness}
\author{fa\_banach\_001, lane 4}
\date{13 August 2026}

\begin{document}
\maketitle

\textbf{Status: literature\_already\_answered.}

\section*{Source question}

In Problem 4 on page 3 of Audenaert and Kittaneh, arXiv:1201.5232v3 \cite{AK}, the authors ask whether, for all positive semidefinite matrices $A,B$, every $p>0$, and every singular-value index $i$,
\begin{equation}\label{eq:question}
 \sigma_i\!\left(A^pB^{1-p}+B^pA^{1-p}\right)
 \leq \sigma_i(A+B).
\end{equation}
Here singular values are arranged in non-increasing order.

\section*{Exact literature answer}

Bottazzi, Elencwajg, Larotonda, and Varela explicitly identify this question in Section 3 of arXiv:1403.7472 \cite{BELV}. Their Example 3.2 gives positive definite $3\times3$ matrices $A,B$ at $p=1/2$ satisfying
\[
 \sigma_3\!\left(A^{1/2}B^{1/2}+B^{1/2}A^{1/2}\right)
 =591.716>443.017=\sigma_3(A+B).
\]
Thus \eqref{eq:question} is false. The supporting authors explicitly cite Audenaert--Kittaneh Problem 4 and state that the example answers it negatively, so this is an exact, provenance-aware literature resolution rather than an inferred implication.

\section*{Supplementary exact $2\times2$ certificate}

The direct attack in this run independently found a smaller symbolic witness. Let
\[
 P=\begin{pmatrix}1&0\\0&0\end{pmatrix},\qquad
 Q=\begin{pmatrix}1/4&\sqrt3/4\\ \sqrt3/4&3/4\end{pmatrix},\qquad
 A=64P,\quad B=Q,\quad p=\frac13.
\]
Both $P$ and $Q$ are orthogonal rank-one projections, hence $A,B\geq0$, $A^{1/3}=4P$, $A^{2/3}=16P$, and $B^{1/3}=B^{2/3}=Q$. Therefore
\[
 L:=A^{1/3}B^{2/3}+B^{1/3}A^{2/3}
 =4PQ+16QP
 =\begin{pmatrix}5&\sqrt3\\4\sqrt3&0\end{pmatrix}.
\]
The eigenvalues of $L^*L$ are $38\pm10\sqrt{13}$. Consequently
\[
 \sigma_2(L)=\sqrt{38-10\sqrt{13}}>1,
\]
because $37>10\sqrt{13}$. Meanwhile
\[
 A+B=\begin{pmatrix}257/4&\sqrt3/4\\ \sqrt3/4&3/4\end{pmatrix},
 \qquad
 \sigma_2(A+B)=\frac{65-\sqrt{4033}}2<1,
\]
where the last inequality follows from $4033>63^2$. Hence already in dimension two,
\[
 \sigma_2\!\left(A^{1/3}B^{2/3}+B^{1/3}A^{2/3}\right)
 >\sigma_2(A+B).
\]
This exact certificate is included as a transparent strengthening of the recorded negative answer; it is not presented as the first resolution of Problem 4.

\section*{Search evidence and scope}

The run searched its lightweight registry, solution, attempt, and proof-gap indexes by arXiv id and core formula keywords, then searched arXiv/general scholarly results by the exact product formula, the source title/authors, and ``singular value'' plus ``counterexample.'' This found arXiv:1403.7472, whose abstract and Example 3.2 directly identify the source problem. No claim is made about the novelty of the supplementary $2\times2$ simplification. The conclusion here concerns only Problem 4; the many other questions collected in arXiv:1201.5232 are outside this packet.

\begin{thebibliography}{9}
\bibitem{AK}
K. M. R. Audenaert and F. Kittaneh,
\emph{Problems and Conjectures in Matrix and Operator Inequalities},
arXiv:1201.5232v3 (2012).

\bibitem{BELV}
T. Bottazzi, R. Elencwajg, G. Larotonda, and A. Varela,
\emph{Inequalities related to Bourin and Heinz means with a complex parameter},
arXiv:1403.7472 (2014), Section 3, Example 3.2.
\end{thebibliography}

\end{document}
